%% kapa: exam=1 round=32 year=2021 period=1 form=A subject=civil_law subject_ko=민법 seq=26 number=26
%% source: Q-Net 감정평가사 기출문제 (cst00309) artlSeq=5208397 file=extracted/1008/5208397/A형/2021년도 제32회 감정평가사 1차 1교시 A형.pdf
%% conversion: kordoc 4.24.51 -> tools/gen1.py (auto); review=proofread-sonnet; flags=none
\begin{question}{26}
구분소유적 공유관계에 관한 설명으로 옳지 않은 것은? (다툼이 있으면 판례에 따름)
\begin{choices}
  \item 구분소유적 공유관계의 해소는 상호명의신탁의 해지에 의한다.
  \item 당사자 내부에 있어서는 각자가 특정매수한 부분은 각자의 단독 소유가 된다.
  \item 구분소유적 공유지분을 매수한 자는 당연히 구분소유적 공유관계를 승계한다.
  \item 제3자의 방해행위가 있으면 공유자는 자기의 구분소유 부분뿐만 아니라 전체 토지에 대하여 공유물의 보존행위로서 그 배제를 구할 수 있다.
  \item 구분소유적 공유관계는 어떤 토지에 관하여 그 위치와 면적을 특정하여 여러 사람이 구분소유하기로 하는 약정이 있어야만 적법하게 성립할 수 있다.
\end{choices}
\end{question}
\answer{3}
